% !TEX root = ../tex/6-5.tex
% 第6章：本文は青、補足は緑
\newcommand{\asmain}{%
    \setbeamercolor{frametitle}{fg=white,bg=main}%
    \setbeamercolor{structure}{fg=main}%
    \setbeamercolor{itemize item}{fg=main}%
    \setbeamercolor{itemize subitem}{fg=main}%
    \setbeamercolor{enumerate item}{fg=main}%
}
\newcommand{\assuppl}{%
    \setbeamercolor{frametitle}{fg=white,bg=examplecol}%
    \setbeamercolor{structure}{fg=examplecol}%
    \setbeamercolor{itemize item}{fg=examplecol}%
    \setbeamercolor{itemize subitem}{fg=examplecol}%
    \setbeamercolor{enumerate item}{fg=examplecol}%
}
% フレームタイトルの「補足」だけを□で囲む
\newcommand{\supplmark}{%
    \setlength{\fboxsep}{1.5pt}%
    \setlength{\fboxrule}{0.7pt}%
    \fbox{\textbf{補足}}%
}
\asmain


%=============================================
\begin{frame}{6.5 Closing Remarks}
    \begin{block}{この章で閉じたこと（論文 6.5）}
        変分・スコア・フローを、\\
        一つの原理的な枠組みに統合する\textbf{要石（keystone）}の章。
    \end{block}
    \begin{itemize}
        \item 共通の「条件付けトリック」
        \item 四つの予測型のオラクル同等性
        \item Fokker--Planck が密度発展の万能法則であること
    \end{itemize}
\end{frame}
%=============================================
\asmain
\begin{frame}{第6章の道筋（全体像）}
    \begin{enumerate}
        \item \textbf{6.1 条件付け}\\
        扱いにくい周辺目標 $\to$ 扱いやすい条件付き目標
        \item \textbf{6.2 学習目標}\\
        スケジュール・予測型・時間重み・サンプリング分布
        \item \textbf{6.3 同等性}\\
        アフィン経路の再パラメータ化と予測型の変換
        \item \textbf{6.4 Fokker--Planck}\\
        三視点を密度レベルでつなぐ万能拘束
    \end{enumerate}
\end{frame}
%=============================================
\asmain
\begin{frame}{二つの核心インサイト}
    \begin{block}{1. 条件付けトリック}
        変分・スコア・フローのいずれも、\\
         intractable な周辺目標を tractable な条件付き目標へ落とす。\\
        これが安定で効率的な学習の「秘伝のタレ」。
    \end{block}
    \begin{block}{2. Fokker--Planck}
        確率密度の時間発展を支配する万能法則。\\
        三視点は、この動力学を尊重する生成過程を、\\
        それぞれ別の言葉で構成している。
    \end{block}
\end{frame}
%=============================================
\asmain
\begin{frame}{四つの予測型は同じ情報を持つ}
    オラクル水準では、次は代数的に変換し合える。
    \begin{center}
        クリーンデータ $x_0$
        \;$\leftrightarrow\;$
        ノイズ $\epsilon$
        \;$\leftrightarrow\;$
        スコア $\nabla\log p_t$
        \;$\leftrightarrow\;$
        速度 $v$
    \end{center}
    \begin{alertblock}{ただし「見た目の差」ではない}
        有限容量ネット・誤差の時間重み・有限ステップ数値解では、\\
        どの量を当てるかで最適化とサンプリングの挙動が変わりうる。
    \end{alertblock}
\end{frame}
%=============================================
\asmain
\begin{frame}{「同等」を区別する}
    章の中心的な教訓は、変換公式の暗記ではなく、\\
    \textbf{どの意味の同等性か}を問うこと。
    \begin{enumerate}
        \item 同じオラクル情報を表す（代数的変換）
        \item 母集団目的が、重みや定数差を除いて同じ
        \item 同じ周辺密度経路 $\{p_t\}$ をたどる
        \item 実務の学習・サンプリングアルゴリズムとして同じ挙動
    \end{enumerate}
    \begin{exampleblock}{注意}
        (1)--(3) が成り立っても、(4) は自動では保証されない。
    \end{exampleblock}
\end{frame}
%=============================================
\asmain
\begin{frame}{設計選択の整理}
    拡散の学習で触る主なレバー：
    \begin{itemize}
        \item 前向きノイズスケジュール $(\alpha_t,\sigma_t)$
        \item 予測ターゲット（$x_0$ / $\epsilon$ / score / velocity）
        \item 時間のサンプリング分布
        \item 時間方向の損失重み
    \end{itemize}
    \begin{block}{見分け方}
        数学的な再パラメータ化で消える差と、\\
        有限モデル・有限計算では残る差を分けて考える。
    \end{block}
\end{frame}
%=============================================
\asmain
\begin{frame}{アフィン経路と密度の共通原理}
    形 $x_t=\alpha_t x_0+\sigma_t\epsilon$ は、時間の再パラメータ化と空間のリスケールで
    \[
        x_t=(1-t)x_0+t\epsilon
    \]
    へ写せる（6.3）。周辺経路の構造は保たれ、時間重みや数値解の都合は変わりうる。
    \begin{block}{密度レベルの標語}
        単純な事前をデータへ運ぶ力学を構成し、\\
        その周辺が所定の確率経路に沿って輸送されるようにする。\\
        FP / 連続の式が、変分・ODE・SDE をつなぐ。
    \end{block}
\end{frame}
%=============================================
\asmain
\begin{frame}{章頭の到達目標との対応}
    章頭の ``After this chapter, you will be able to...'' への対応。
    \begin{itemize}
        \item 条件付け原理（周辺 $\to$ 条件付き）を三視点で説明できる
        \item 四つの予測型を変換し、オラクル情報の同一性を言える
        \item 設計レバーのうち、再パラメータ化と実務差を切り分けられる
        \item 「同等」の複数の意味を区別できる
        \item アフィン経路の時間・空間変換の保存／非保存を追える
        \item 連続の式 / FP が三記述をつなぐ密度原理だと説明できる
    \end{itemize}
\end{frame}
%=============================================
\assuppl
\begin{frame}{\supplmark：到達目標チェックリスト}
    \footnotesize
    \begin{enumerate}
        \item 変分・スコア・フローで、同じ条件付けパターンを指摘できるか
        \item $x_0$/$\epsilon$/score/$v$ のオラクル変換を、非退化スケジュールで書けるか
        \item スケジュール・予測型・時間分布・重みの役割を分けて言えるか
        \item 「情報同等」「目的同等」「周辺同等」「アルゴリズム同等」を区別できるか
        \item $\alpha_t,\sigma_t$ 経路を線形スケジュールへ写すとき、何が変わるか説明できるか
        \item Continuity / FP を、三視点の共通拘束として一文で言えるか
    \end{enumerate}
\end{frame}
%=============================================
\asmain
\begin{frame}{Reader route（初読の指針）}
    この章は本全体の概念的な要石である。初読では次に集中する。
    \begin{enumerate}
        \item 条件付け原理
        \item 四つの予測型の関係
        \item 連続の式 / Fokker--Planck による三視点の接続
    \end{enumerate}
    \begin{block}{中心的な教訓}
        変換公式を暗記することではなく、\\
        「どの同等性か／何が保たれるか」を問う習慣である。
    \end{block}
    \vfill
    \footnotesize
    重み・スケジュール・パラメータ化の細部は、後から実務差の理解のために再訪する。
\end{frame}
%=============================================
\asmain
\begin{frame}{核心の一文}
    \begin{block}{第6章の核心}
        拡散の三視点は、同じ確率経路の周辺を運ぶ生成子の別表現である。\\
        学習は条件付き目標への回帰、密度の法則は Fokker--Planck である。
    \end{block}
    \begin{itemize}
        \item オラクルでは予測型もスケジュールも大きく変換できる
        \item 実務では有限性のため、表現の選びが効きうる
        \item 生成 $=$ 微分方程式を解く、という土台が次パートへ続く
    \end{itemize}
\end{frame}
%=============================================
\asmain
\begin{frame}{この先の地図（Part C / D）}
    「何を・なぜ」を統一したので、次は「どう使う・どう速くする」。
    \begin{enumerate}
        \item \textbf{Part C（推論時）}\\
        軌跡の制御（第8章）、高速な数値ソルバ（第9章）
        \item \textbf{Part D（少数ステップ）}\\
        蒸留（第10章）、スクラッチの高速生成器（第11章）
    \end{enumerate}
    \begin{exampleblock}{共通の足場}
        生成 $=$ 微分方程式を解くこと。\\
        統一理解がその制御と加速のプラットフォームになる。
    \end{exampleblock}
\end{frame}
%=============================================
\asmain
\begin{frame}{第7章へ：Optimal Transport}
    分布から分布への写像（生成はその特別な場合）が次の主題。
    \begin{itemize}
        \item Flow Matching：ソース$\to$ターゲットの速度場学習
        \item 古典 OT / エントロピー OT / Schr\"odinger bridge
        \item 確率的視点（SDE）はエントロピー OT と親和的
        \item PF-ODE は既知の輸送目的と一般には一致しない
    \end{itemize}
    \begin{block}{つなぎ}
        第6章の「同じ周辺・異なる力学」が、\\
        「どの輸送を実現しているか」という第7章の問いに直結する。
    \end{block}
\end{frame}
%=============================================
\asmain
\begin{frame}{第6章のまとめ}
    \begin{itemize}
        \item 条件付けが、三視点に共通する学習の仕組みである
        \item 四予測型はオラクル同等だが、実務挙動は異なりうる
        \item 設計選択は、再パラメータ化と真の自由度を分けて読む
        \item Continuity / FP が変分・ODE・SDE を密度で統一する
        \item 次は第7章：Optimal Transport との関係を問う
    \end{itemize}
\end{frame}
